\documentclass[11pt]{article}

\usepackage[a4paper,margin=1in]{geometry}
\usepackage{amsmath,amssymb,amsthm}
\usepackage{graphicx}
\usepackage[hidelinks]{hyperref}

\newtheorem{theorem}{Theorem}
\newtheorem{lemma}{Lemma}
\newtheorem{corollary}{Corollary}

\newcommand{\eps}{\varepsilon}
\newcommand{\norm}[1]{\left\lVert #1 \right\rVert}
\newcommand{\ball}{B}
\newcommand{\id}{I}

\title{Weak-star Delta-points Do Not Appear In Duals Of Subspaces Of Small Unconditional Basis Spaces}
\author{Automatic math-discovery smoke test}
\date{June 14, 2026}

\begin{document}
\maketitle

\section*{Source Problem}

The source paper is Trond A. Abrahamsen, Ramon J. Aliaga, Vegard Lima,
Andre Martiny, Yoel Perreau, Antonin Prochazka, and Triinu Veeorg,
\emph{Delta-points and their implications for the geometry of Banach spaces},
arXiv:2303.00511 \cite{AALMPPV2023}.  The relevant question appears as
Question 5.18 on page 23 of the arXiv PDF, after Corollary 5.17.

\begin{center}
\includegraphics[width=\linewidth]{figures/open_problem_crop.png}
\end{center}

In text form, the question is:
\begin{quote}
Can we conclude that \(X^*\) contains no weak* \(\Delta\)-points if \(X\) is
a subspace of a Banach space with a shrinking \(k\)-unconditional basis for
\(k<2\)?
\end{quote}

\section*{Result Type}

This packet gives a \textbf{candidate full affirmative answer} to Question
5.18.  It proves the statement as asked, with the weak-star topology on
\(X^*\) coming from the predual \(X\).

\section*{Definitions And Source Inputs}

For a Banach space \(E\), write \(\ball_E\) for its closed unit ball.  A point
\(x^*\in S_{E^*}\) is a weak-star \(\Delta\)-point if every weak-star slice of
\(\ball_{E^*}\) containing \(x^*\) has points arbitrarily close to distance
\(2\) from \(x^*\).  It is a weak-star super \(\Delta\)-point if the same
condition holds for every relatively weak-star open subset of \(\ball_{E^*}\)
containing \(x^*\).

We use the following theorem from the source paper \cite[Theorem 5.6]{AALMPPV2023}:
if \(E^*\) contains a weak-star \(\Delta\)-point, then \(E^*\) contains a
weak-star super \(\Delta\)-point.

We will also use the standard fact that if \(T:E\to E\) is compact, then
\(T^*:(\ball_{E^*},w^*)\to(E^*,\norm{\cdot})\) is continuous.  This is the
compact-adjoint continuity used in the proof of \cite[Theorem 5.12]{AALMPPV2023}.

\section*{Proof Intuition}

The source paper already proves a no-weak-star-\(\Delta\) theorem when the
space itself has a shrinking \(k\)-unconditional basis.  The only obstacle in
Question 5.18 is that a subspace \(X\subset Y\) need not inherit the finite-rank
coordinate projections of \(Y\).  Passing to the dual removes this obstacle:
\(X^*\) is the quotient of \(Y^*\) by restriction.  A weak-star convergent net
in the quotient unit ball can be lifted to the ambient dual unit ball, and
weak-star compactness gives a convergent subnet.  The ambient shrinking
unconditional basis supplies a uniform distance bound strictly below \(2\) on
that subnet.  Since quotient distance is no larger than ambient distance, the
same obstruction holds in \(X^*\).

\section*{A Quotient Stability Lemma}

\begin{lemma}\label{lem:ambient-bound}
Let \(E\) be a Banach space.  Suppose there is a family
\((T_\lambda)_{\lambda\in\Lambda}\) of compact operators on \(E\) such that
\[
  r:=\sup_{\lambda\in\Lambda}\norm{\id_E-T_\lambda}<2
  \quad\text{and}\quad
  T_\lambda^*x^*\to x^*
  \quad\text{in norm for every }x^*\in E^*.
\]
Then there is a constant \(c<2\) such that every weak-star convergent net
\((x_\alpha^*)\subset \ball_{E^*}\), say \(x_\alpha^*\to x^*\), has the
following cofinal closeness property: for every index \(\alpha_0\), there is
\(\alpha\succeq\alpha_0\) with
\[
  \norm{x_\alpha^*-x^*}\le c .
\]
\end{lemma}

\begin{proof}
Choose \(\eps>0\) so that \(c:=r+2\eps<2\).  Fix a weak-star convergent net
\((x_\alpha^*)\subset \ball_{E^*}\), with limit \(x^*\), and fix an index
\(\alpha_0\).  Pick \(\lambda\) such that
\[
  \norm{x^*-T_\lambda^*x^*}\le \eps.
\]
Since \(T_\lambda\) is compact, \(T_\lambda^*x_\alpha^*\to T_\lambda^*x^*\)
in norm.  Hence there is an index \(\beta\) such that
\(\norm{T_\lambda^*(x_\alpha^*-x^*)}\le\eps\) whenever
\(\alpha\succeq\beta\).  Choose \(\alpha\succeq\alpha_0,\beta\).  Then
\[
\begin{aligned}
\norm{x_\alpha^*-x^*}
&\le
\norm{x^*-T_\lambda^*x^*}
+\norm{T_\lambda^*(x_\alpha^*-x^*)}
+\norm{x_\alpha^*-T_\lambda^*x_\alpha^*}  \\
&\le \eps+\eps+\norm{\id_E-T_\lambda}
\le r+2\eps
=c<2 .
\end{aligned}
\]
\end{proof}

\begin{lemma}\label{lem:quotient-transfer}
Let \(X\) be a closed subspace of a Banach space \(Y\).  Assume that \(Y\)
satisfies the conclusion of Lemma \ref{lem:ambient-bound}; that is, there is
\(c<2\) such that every weak-star convergent net in \(\ball_{Y^*}\) has a
cofinal point within distance \(c\) of its limit.  Then \(X^*\) has no
weak-star super \(\Delta\)-points.
\end{lemma}

\begin{proof}
Let \(Q:Y^*\to X^*\) be the restriction map.  It is weak-star to weak-star
continuous and, by Hahn--Banach, \(Q(\ball_{Y^*})=\ball_{X^*}\).

First we prove the same cofinal closeness property in \(\ball_{X^*}\).
Let \((u_\alpha^*)\subset\ball_{X^*}\) be a weak-star convergent net with
limit \(u^*\), and fix \(\alpha_0\).  Restrict to the tail
\(\{\alpha:\alpha\succeq\alpha_0\}\).  For each \(\alpha\) in this tail,
choose \(v_\alpha^*\in\ball_{Y^*}\) with \(Qv_\alpha^*=u_\alpha^*\).  By
weak-star compactness of \(\ball_{Y^*}\), the net \((v_\alpha^*)\) has a
weak-star convergent subnet \((v_\beta^*)\), with limit \(v^*\in\ball_{Y^*}\).
Continuity of \(Q\) gives \(Qv^*=u^*\).  By the assumed ambient property,
there is a subnet index \(\beta\) such that
\[
  \norm{v_\beta^*-v^*}\le c.
\]
Therefore
\[
  \norm{u_\beta^*-u^*}
  =\norm{Q(v_\beta^*-v^*)}
  \le \norm{v_\beta^*-v^*}
  \le c.
\]
This proves the cofinal closeness property in \(\ball_{X^*}\).

Now suppose, for contradiction, that \(u^*\in S_{X^*}\) is a weak-star super
\(\Delta\)-point.  Direct the set of pairs \((U,n)\), where \(U\) is a
relatively weak-star open neighborhood of \(u^*\) in \(\ball_{X^*}\) and
\(n\in\mathbb{N}\), by reverse inclusion in \(U\) and increasing \(n\).  Since
\(u^*\) is weak-star super \(\Delta\), for every \((U,n)\) we may choose
\(u_{U,n}^*\in U\) with
\[
  \norm{u_{U,n}^*-u^*}>2-\frac1n .
\]
This net converges weak-star to \(u^*\), and eventually its distance from
\(u^*\) is larger than \(c\), because \(c<2\).  This contradicts the cofinal
closeness property just proved.  Thus \(X^*\) has no weak-star super
\(\Delta\)-points.
\end{proof}

\section*{Main Result}

\begin{theorem}\label{thm:main}
Let \(X\) be a closed subspace of a Banach space \(Y\).  If \(Y\) has a
shrinking \(k\)-unconditional basis for some \(k<2\), then \(X^*\), equipped
with the weak-star topology induced by \(X\), contains no weak-star
\(\Delta\)-points.
\end{theorem}

\begin{proof}
Let \((e_n)\) be the shrinking \(k\)-unconditional basis of \(Y\), and let
\(P_n:Y\to Y\) be the \(n\)-th partial sum projection.  The operators \(P_n\)
are finite-rank, hence compact.  Since the basis is shrinking, \(P_n^*y^*\to
y^*\) in norm for every \(y^*\in Y^*\).  Since the basis is
\(k\)-unconditional, the tail projections satisfy
\[
  \norm{\id_Y-P_n}\le k
  \quad\text{for every }n.
\]
Thus Lemma \ref{lem:ambient-bound} applies to \(E=Y\), with \(r\le k<2\).
Lemma \ref{lem:quotient-transfer} then shows that \(X^*\) contains no
weak-star super \(\Delta\)-points.

Finally, by \cite[Theorem 5.6]{AALMPPV2023}, the existence of a weak-star
\(\Delta\)-point in \(X^*\) would imply the existence of a weak-star super
\(\Delta\)-point in \(X^*\).  This is impossible by the previous paragraph.
Therefore \(X^*\) contains no weak-star \(\Delta\)-points.
\end{proof}

\begin{corollary}
Question 5.18 of Abrahamsen--Aliaga--Lima--Martiny--Perreau--Prochazka--Veeorg
has an affirmative answer.
\end{corollary}

\section*{Verification Notes}

The main point to verify is Lemma \ref{lem:quotient-transfer}.  The lift
\((v_\alpha^*)\) is taken only on an arbitrary tail of the original net in
\(\ball_{X^*}\), so the weak-star compact subnet still corresponds to indices
cofinal beyond the prescribed starting point.  Since the restriction map
\(Q:Y^*\to X^*\) is weak-star continuous, the lifted subnet limit restricts to
the original quotient limit.  The quotient norm inequality
\(\norm{Qz^*}\le\norm{z^*}\) then transfers the strict \(<2\) distance bound.

No numerical computation is used.  The proof is independent of separability
except for the standard countable-basis setting in the source question.

\section*{Novelty And Literature Check}

Before promotion, the run indexes
\[
\texttt{registry\_index.tsv},\quad
\texttt{solutions/index.tsv},\quad
\texttt{attempts/index.tsv},\quad
\texttt{proof\_gaps/index.tsv}
\]
were searched for arXiv:2303.00511, the title, and core phrases including
\(\Delta\)-points, Daugavet points, weak-star \(\Delta\)-points, subspaces,
and shrinking \(k\)-unconditional bases.  No duplicate packet, prior attempt,
or proof-gap entry was found.

A bounded web search on June 14, 2026 for exact and close variants of Question
5.18, including phrases such as ``weak* Delta-points'', ``subspace'',
``shrinking k-unconditional basis'', and ``2303.00511'', found the source paper
but no separate later paper explicitly resolving this question.  A local corpus
search across the parsed arXiv sources and run target-pool indexes likewise
found the source question and background mentions of the source paper, but no
explicit later answer to Question 5.18.

\section*{Human Review Recommendation}

Send for human review as a candidate full solution.  The recommended review
focus is the subnet/quotient step in Lemma \ref{lem:quotient-transfer}; once
that lemma is accepted, the specialization to shrinking \(k\)-unconditional
bases is exactly the compact-approximation mechanism used in the source paper.

\begin{thebibliography}{9}

\bibitem{AALMPPV2023}
T.~A. Abrahamsen, R.~J. Aliaga, V.~Lima, A.~Martiny, Y.~Perreau,
A.~Prochazka, and T.~Veeorg,
\emph{Delta-points and their implications for the geometry of Banach spaces},
arXiv:2303.00511, 2023.

\bibitem{BowersKalton2014}
A.~Bowers and N.~J. Kalton,
\emph{An introductory course in functional analysis},
Universitext, Springer, 2014.

\end{thebibliography}

\end{document}
